\chapter{Testing}
\label{chap:testing}

\section{Strategia di test}
\label{sec:test_strategy}

La strategia di test si articola su tre livelli: test unitari, test di integrazione
e test di carico. L'obiettivo \`e verificare sia la correttezza funzionale dei singoli
componenti sia il comportamento complessivo del sistema sotto carico realistico.

\section{Test unitari del backend}
\label{sec:unit_backend}

I test unitari del backend si trovano in \texttt{backend/tests/} e coprono:

\begin{itemize}[leftmargin=*]
    \item \texttt{test\_upload.py}: flusso two-phase di upload (creazione documento,
          generazione presigned URL, conferma con \texttt{size\_bytes}).
    \item \texttt{test\_s3\_exist.py}: verifica dell'esistenza del bucket S3.
    \item \texttt{test\_s3\_copy.py}: operazione di copia oggetto su SeaweedFS.
    \item \texttt{test\_s3\_missing.py}: gestione corretta degli errori per oggetti
          inesistenti.
\end{itemize}

\section{Test unitari dell'agent worker}
\label{sec:unit_agent}

Il file \texttt{agent\_worker/test\_agent.py} contiene test unitari per il filing
agent, verificando:

\begin{itemize}[leftmargin=*]
    \item La correttezza della struttura di output (\texttt{FilingOutput} conforme
          allo schema Pydantic).
    \item Il comportamento dell'agente in assenza di cartelle (\texttt{folder\_id = null}).
    \item La propagazione del reasoning testuale nella risposta.
\end{itemize}

\TODO{Aggiungere test con mock del client LLM per testare il behavior dell'agente
senza consumare token reali; aggiungere test per la logica di confidence threshold;
aggiungere test per la validazione UUID del folder\_id.}

\section{Test di integrazione (script bash)}
\label{sec:integration_tests}

Gli script di integrazione in \texttt{backend/tests/} eseguono scenari end-to-end
contro il sistema in esecuzione:

\begin{description}[leftmargin=*]
    \item[\texttt{test.sh}] Happy path completo: registrazione utente, login, creazione
          cartella, upload documento, listing, download, cancellazione.
    \item[\texttt{test\_scenario.sh}] Scenario multi-utente: due utenti dello stesso
          dominio, condivisione di una cartella, accesso con permesso VIEWER e EDITOR,
          tentativo di accesso cross-user non autorizzato.
    \item[\texttt{test\_agent\_scenario.sh}] Flusso email $\rightarrow$ classificazione
          $\rightarrow$ proposta inbox: configura un account IMAP di test, verifica la
          creazione del job, attende l'elaborazione e controlla la presenza della proposta
          nell'\emph{inbox view}.
\end{description}

\section{Test di carico}
\label{sec:stress_test}

Lo script \texttt{backend/tests/stress-test.sh} esegue richieste concorrenti verso gli
endpoint pi\`u utilizzati, misurando:

\begin{itemize}[leftmargin=*]
    \item Throughput (richieste/secondo) per \texttt{GET /documents} e \texttt{GET /folders}.
    \item Latenza P50, P95, P99 sotto carico sostenuto.
    \item Comportamento del connection pool PostgreSQL: assenza di errori
          ``connection refused'' con \texttt{max\_overflow = 10}.
    \item Assenza di memory leak nel backend sotto stress prolungato.
\end{itemize}

\TODO{Sostituire lo stress test bash con uno strumento dedicato (Locust o k6) per
report pi\`u dettagliati; aggiungere test di carico per la pipeline agentica (job
concorrenti); misurare il throughput massimo sostenibile prima di saturazione del pool DB.}

\section{Test di failure injection}
\label{sec:failure_injection}

I test di failure injection verificano il comportamento del sistema sotto condizioni
degradate. Sono eseguibili manualmente tramite comandi Docker Compose e sono documentati
nel \Cref{chap:failure}.

\TODO{Automatizzare i test di failure injection con Toxiproxy o Chaos Monkey per
simulare: latenza di rete, packet loss, crash di singoli container, saturazione del
connection pool. Integrare nella pipeline CI come job separato.}

\section{Risultati}
\label{sec:test_results}

\TODO{Inserire i risultati quantitativi dei test di carico: latenza misurata P50/P95/P99,
throughput massimo, tasso di errori sotto carico. Includere grafici Grafana o output
di Locust.}
